\documentclass[12pt]{article}
\usepackage{fullpage}
\usepackage{amsmath,amssymb,amsthm}

\newtheorem{lemma}{Lemma}
\newcommand{\ol}{\overline}

\begin{document}

\begin{center}
{\bf A counterexample}
\end{center}

The answer is no.  Already in the free group \(G=\langle a,b,c\rangle\) there is a counterexample.  Write \(\ol x=x^{-1}\), and consider
\[
        w=\ol a\,a\, b\,\ol b\,\ol a\,a\, c\,\ol c\,\ol a\,a .
\]
This word is freely trivial (delete the five displayed adjacent inverse
pairs), and it is quasi-reduced: after each adjacent inverse pair the
next letter, if present, is not the first letter of that pair.

We shall use two elementary facts about the definition.  Compactify the
upper half-plane to a disk, and put the marked boundary points in cyclic
order.  For a matching \(M\), let \(\Gamma(M)\) be the graph obtained
from \(M\) by making every crossing a vertex.

\begin{lemma}
For every matching satisfying the conditions in the question, the graph
\(\Gamma(M)\) is a forest.
\end{lemma}

\begin{proof}
If \(\Gamma(M)\) contained a cycle, then, since the boundary vertices
have valence one, that cycle would lie in the interior of the disk.
Choose an innermost such cycle.  It bounds a complementary region whose
intersection with \(M\) has a component with no boundary endpoint,
contrary to the requirement that every such component have two boundary
endpoints.
\end{proof}

\begin{lemma}
Suppose that a letter \(x\) occurs in \(w\) exactly once and \(x^{-1}\)
occurs exactly once.  Then, in every matching, those two endpoints are
joined by an uncrossed interval.
\end{lemma}

\begin{proof}
Let the two endpoints be \(p\) and \(q\).  At \(p\), the two local sides
of the incident edge belong to two face-boundary components.  By the
label condition, both of these components must end at \(q\), because
there is no other endpoint labelled by the inverse of the label at
\(p\).  By the previous lemma, the component of \(\Gamma(M)\) containing
\(p\) is a tree, so there is a unique path from \(p\) to \(q\).

If that path met a crossing, take the first crossing encountered from
\(p\).  The two local sides of the incoming edge leave that crossing
along two different half-edges, hence into two different components of
the tree with the crossing removed.  At most one of those components
contains \(q\), a contradiction.  Therefore the path from \(p\) to \(q\)
has no crossing, and is a single uncrossed interval.
\end{proof}

In our word the letters \(b,\ol b,c,\ol c\) each occur only in the
displayed adjacent inverse pairs.  Hence every matching contains the
uncrossed intervals \((3,4)\) and \((7,8)\).  Since their endpoints are
adjacent, these intervals cut off empty boundary disks and are
boundary-parallel.  Deleting
them and renumbering the remaining endpoints
\[
        1,2,5,6,9,10
\]
gives a cell-preserving isomorphism from \(F_w\) to the analogous
complex for the shorter word (using the same definition)
\[
        u=\ol a\,a\,\ol a\,a\,\ol a\,a .
\]
Thus it remains only to compute this six-endpoint complex.

For \(u\), odd endpoints have label \(\ol a\) and even endpoints have
label \(a\).  We use the notation
\[
        ij|kl|mn
\]
for the endpoint pairing \(\{i,j\},\{k,l\},\{m,n\}\).  By the first
lemma, with three intervals no two intervals can cross twice and not
all three pairs can cross.  Hence an isotopy class is determined by its
endpoint pairing, with a crossing exactly when two pairs interleave.

Conversely, if the interleaving graph of a three-pair chord diagram is a
forest, then its face-boundary components are arcs joining boundary
endpoints.  Pushing such a component slightly into the adjacent face
gives an arc disjoint from the diagram; the two components of the disk
cut along this arc must each contain an even number of marked endpoints:
an interval with endpoints on different sides would have to cross the
pushed-off arc.  Equivalently, the number of marked endpoints strictly
between the two ends along either boundary arc is even.  Thus the two
ends of the face-boundary component have opposite parity.
For the alternating word \(u\), this is exactly the label condition.
Therefore all endpoint pairings except the unique all-crossing pairing
\(14|25|36\) occur.

The cells are as follows.  The five vertices are
\[
\begin{array}{lll}
v_0=12|34|56, & v_1=12|36|45, & v_2=14|23|56,\\
v_3=16|23|45, & v_4=16|25|34. &
\end{array}
\]
The six edges are
\[
\begin{array}{lll}
e_{01}=12|35|46, & e_{02}=13|24|56, & e_{13}=13|26|45,\\
e_{23}=15|23|46, & e_{04}=15|26|34, & e_{34}=16|24|35,
\end{array}
\]
where the subscripts indicate the two endpoint vertices.  The three
squares are
\[
\begin{array}{lll}
Q_{12}=13|25|46, & Q_{14}=14|26|35, & Q_{24}=15|24|36,
\end{array}
\]
with boundaries
\[
\begin{aligned}
\partial Q_{12}&=e_{01}\cup e_{13}\cup e_{23}\cup e_{02},\\
\partial Q_{14}&=e_{01}\cup e_{13}\cup e_{34}\cup e_{04},\\
\partial Q_{24}&=e_{02}\cup e_{23}\cup e_{34}\cup e_{04}.
\end{aligned}
\]
These boundary relations are exactly the two local resolutions of each
crossing.

Thus the one-skeleton is the theta graph with poles \(v_0,v_3\) and
three length-two paths through \(v_1,v_2,v_4\).  The three square cells
fill the three cycles obtained by choosing two of those paths.  This is
the standard cell structure on \(S^2\) obtained from three meridional
lunes.  Hence
\[
        F_w \cong S^2.
\]
In particular \(F_w\) is not contractible.

\end{document}
